\documentclass[11pt]{article}

\usepackage[margin=1in]{geometry}
\usepackage{amsmath,amssymb,amsthm}
\usepackage{enumitem}
\usepackage{hyperref}

\newtheorem{theorem}{Theorem}
\newtheorem{lemma}{Lemma}
\newtheorem{remark}{Remark}

\newcommand{\Mg}{\mathcal M_g}
\newcommand{\Mbar}{\overline{\mathcal M}}
\newcommand{\T}{\mathcal T}
\newcommand{\DeltaG}{\Delta_\Gamma}

\begin{document}

\title{The proposed non-variation statement is false as stated}
\author{}
\date{}
\maketitle

\begin{abstract}
We show that the statement in the question cannot be proved from the
stated hypotheses: the hypotheses allow a local holomorphic motion of the
constant boundary factor.  The construction uses a standard totally
geodesic Teichmuller surface coming from cyclic covers and then translates
the constant factor in a plumbing neighbourhood of a boundary component.
\end{abstract}

\section*{Counterexample}

We work in the usual orbifold sense for moduli spaces.  Equivalently, one
may pass throughout to a sufficiently high level cover to remove stabilizers;
the construction is local in an orbifold chart.

\begin{lemma}[A totally geodesic surface with a constant boundary factor]
There is a complex two-dimensional totally geodesic suborbifold
\(N\subset \mathcal M_6\) whose Deligne--Mumford closure has a boundary
component
\[
 B\subset \Delta_\Gamma \cong \mathcal M_{2,1}\times \mathcal M_{4,1}
\]
of the form
\[
 B=\{Y_0\}\times C,
\]
where \(Y_0\in \mathcal M_{2,1}\) is fixed and \(C\subset \mathcal M_{4,1}\)
is a non-constant curve.
\end{lemma}

\begin{proof}
Let \(N\) be the locus of cyclic degree \(5\) covers of \(\mathbb P^1\)
branched at five ordered points with local monodromy \(1\) at each branch
point; equivalently, on an affine chart,
\[
 y^5=\prod_{j=1}^5 (x-z_j).
\]
By Riemann--Hurwitz,
\[
 2g-2=5(-2)+5(5-1)=10,
\]
so \(g=6\).  Varying the five branch points modulo
\(\operatorname{PGL}_2\) gives a two-dimensional locus.  This locus is the
fixed locus of an order \(5\) mapping class whose quotient orbifold is a
sphere with five cone points of order \(5\).  The fixed locus in
Teichmuller space is identified with the Teichmuller space of that quotient
orbifold.  Since Teichmuller geodesics between two points are unique and
the order \(5\) mapping class acts by Teichmuller isometries, this fixed
locus is totally geodesic.  Hence its image \(N\subset \mathcal M_6\) is a
totally geodesic suborbifold.

Now degenerate the base \(\mathbb P^1\) to a two-component stable rational
curve, with two branch points on one component and three on the other.  The
corresponding admissible cover has two components meeting in one node.  On
the component carrying two branch points, the node has monodromy
\(-2\equiv 3\pmod 5\), so the cover is a cyclic degree \(5\) cover of
\(\mathbb P^1\) branched at three points with monodromies \((1,1,3)\).
Again by Riemann--Hurwitz,
\[
 2g_1-2=5(-2)+3(5-1)=2,
\]
so \(g_1=2\).  Since a three-pointed rational curve has no moduli, this
genus \(2\) pointed component is a fixed point
\(Y_0\in \mathcal M_{2,1}\).

On the other component, the cover is branched at four points, with
monodromies \((1,1,1,2)\).  Its genus is \(4\), since
\[
 2g_2-2=5(-2)+4(5-1)=6.
\]
The cross-ratio of the four special points varies, producing a curve
\(C\subset \mathcal M_{4,1}\).  Thus the relevant boundary component is
\[
 B=\{Y_0\}\times C\subset \mathcal M_{2,1}\times \mathcal M_{4,1}.
\]
\end{proof}

We now use the standard plumbing description of a neighbourhood of a
boundary stratum.  After shrinking around the closure of \(B\), there is a
stratified analytic neighbourhood \(W\subset \overline{\mathcal M}_6\) of
\(\overline B\) of the form
\[
 W\simeq U\times W',
\]
where \(U\subset \mathcal M_{2,1}\) is a small coordinate neighbourhood of
\(Y_0\), and \(W'\) contains the \(\mathcal M_{4,1}\)-factor together with
the relevant plumbing parameters.  The Deligne--Mumford strata in \(W\) are
products of \(U\) with strata in \(W'\).

Choose a non-constant holomorphic disc
\[
 t\longmapsto Y_t\in U,\qquad Y_0=Y_{t=0}.
\]
Let \(\tau_t:U\to U\) be a holomorphic local motion with
\(\tau_0=\mathrm{id}\) and \(\tau_t(Y_0)=Y_t\).  Define
\[
 \Phi_t:W\to W,\qquad \Phi_t(u,w')=(\tau_t(u),w').
\]
Thus \(\Phi_t\) preserves every boundary stratum and moves only the
\(\mathcal M_{2,1}\)-factor.

Let \(V_0\) be a sufficiently small collar neighbourhood of \(B\) inside
the totally geodesic surface \(N\cap \mathcal M_6\).  Then \(V_0\) is a
smooth complex two-dimensional submanifold of \(\mathcal M_6\), and it is
totally geodesic because it is an open subset of \(N\).  For \(t\) small,
set
\[
 V_t:=\Phi_t(V_0)\cap \mathcal M_6.
\]
Equivalently,
\[
 \mathcal V:=\{(\Phi_t(X),t):X\in V_0,\ t\in \Delta\}
 \subset \mathcal M_6\times \Delta
\]
is a holomorphic family of smooth two-dimensional submanifolds.

Since \(\Phi_t\) preserves the Deligne--Mumford stratification, for every
stratum \(\Delta_G\) one has
\[
 \overline{V_t}\cap \Delta_G
 =
 \Phi_t\bigl(\overline{V_0}\cap \Delta_G\bigr).
\]
Hence the boundary intersections form smoothly trivial families.  Moreover,
because \(\Phi_t\) preserves strata and commutes with taking closures inside
the chosen product plumbing neighbourhood, the compatibility condition in
the question is also preserved.

For the particular stratum
\[
 \Delta_\Gamma\cong \mathcal M_{2,1}\times \mathcal M_{4,1},
\]
we obtain
\[
 \partial V_t
 =
 \overline{V_t}\cap \Delta_\Gamma
 =
 \Phi_t(B)
 =
 \{Y_t\}\times C.
\]
Therefore
\[
 \pi_1(\partial V_t)=\{Y_t\}.
\]
By construction \(Y_t\neq Y_0\) for \(t\neq 0\).  Thus the conclusion
claimed in the question,
\[
 \pi_1(\partial V_t)=\{Y_0\}\quad\text{for all }t,
\]
fails.

\begin{theorem}
The proposed non-variation statement is false under the hypotheses as
stated.
\end{theorem}

\begin{proof}
The family \(V_t\) constructed above satisfies all stated hypotheses: \(V_0\)
is a smooth complex two-dimensional totally geodesic submanifold, the
boundary intersections with every stratum are smoothly trivial and
compatible, and \(\partial V_t\subset \Delta_\Gamma\) is one-dimensional and
constant on the \(\mathcal M_{2,1}\)-factor.  However that constant factor
is \(Y_t\), which varies holomorphically with \(t\).  Hence the asserted
conclusion is false.
\end{proof}

\begin{remark}
The failure comes from a simple local phenomenon: in plumbing coordinates
near a boundary stratum, a factor on which the central boundary curve is
constant can be moved independently while leaving all incidence relations
among boundary strata unchanged.  Therefore additional hypotheses, such as
a global rigidity condition on the deformation, would be needed to obtain
the desired non-variation conclusion.
\end{remark}

\begin{thebibliography}{9}

\bibitem{Bers}
L.~Bers,
\emph{Spaces of degenerating Riemann surfaces},
Annals of Mathematics Studies, vol.~79, Princeton Univ. Press, 1974,
pp.~43--55.

\bibitem{DM}
P.~Deligne and D.~Mumford,
\emph{The irreducibility of the space of curves of given genus},
Inst. Hautes \'{E}tudes Sci. Publ. Math. \textbf{36} (1969), 75--109.

\bibitem{Gardiner}
F.~P. Gardiner,
\emph{Teichmuller Theory and Quadratic Differentials},
Pure and Applied Mathematics, Wiley, 1987.

\bibitem{Harvey}
W.~J. Harvey,
\emph{On branch loci in Teichmuller space},
Trans. Amer. Math. Soc. \textbf{153} (1971), 387--399.

\bibitem{HK}
J.~Hubbard and S.~Koch,
\emph{An analytic construction of the Deligne--Mumford compactification of
the moduli space of curves},
J. Differential Geom. \textbf{98} (2014), no.~2, 261--313.

\end{thebibliography}

\end{document}